\begin{table}[H]
\centering
\caption*{\textbf{Quadro 1 - Decisões centrais da MOSAIC 0.2}}
\small
\renewcommand{\arraystretch}{1.18}
\begin{tabularx}{\linewidth}{|>{\bfseries\raggedright\arraybackslash}p{0.31\linewidth}|Y|}
\hline
Decisão & Efeito arquitetural \\
\hline
Planejar por objetivos orientados a resultados &
A granularidade do plano não copia tools nem arquivos de skills. \\
\hline
Tratar skills como instruções e tools como operações &
O modelo continua responsável por interpretação, transformação e decisão. \\
\hline
Selecionar bundles ordenados de zero a várias skills &
Um objetivo combina comportamentos complementares sem virar vários nós. \\
\hline
Formar o menu por \(T_{\mathrm{base}}\) e pelas tools do bundle &
Um bundle vazio pode usar operações gerais sem exigir uma skill artificial. \\
\hline
Capturar observações no runtime &
Cada retorno de tool executável vira evidência ordenada; critérios citam índices
locais e a decisão não reproduz identificadores do provedor. \\
\hline
Associar revisão ao conjunto completo de evidências &
Uma revisão localizada recebe todas as observações do nó e preserva os
resultados já concluídos. \\
\hline
\end{tabularx}
\par\smallskip\raggedright\small Fonte: elaboração própria (2026).
\end{table}
